\section{Conclusion}\label{sec:concl}
% =============================================================================================
% Owner: front writer. Plan: OUTLINE.md, Section 10. Every number is taken from a generated or
% condensed table that the result sections input: tab:e3-power (manual_e3_power_main),
% tab:e3-pooled-tests, tab:e3-verdicts, tab:e1-size-sim (manual_e1_size_sim_main),
% tab:e2-relative-boundary, tab:e2-power, tab:e2dep-power, tab:e2-llr, tab:e3-perticker-null,
% tab:e4-correlation, tab:e4-utility, tab:e4-utility-agg, tab:e4-tests, tab:e4-link,
% tab:manual-e5-sensitivity, tab:manual-e5-split, tab:e5-simpower, tab:e5-arch (e5_arch_sim),
% tab:e5-arch-real (manual_e5_arch_real), tab:manual-theory-calibration (size-adjusted power),
% tab:manual-e1-sim-stat (correlation of T with its HAC s.e.), or from the text of
% sec:emp-witness (0.65) and sec:emp-likelihood (1/6). The configuration counts of the E5 grid
% (13, 11 and 9 of 14 BAD; 69 of 70 absolute rejections) are counted from
% manual_e5_sensitivity.tex. Two derived counts were recomputed with uv run python from the raw
% results (revision v1, 2026-11-16): (i) alternative split, design D2
% (results/e5/e5_split_power_raw.csv, pmed = min(1, 2 median) over the five seeds per
% replication): the alternative and the boundary are rejected in the same 15 of 25 replications;
% every rejected alternative has pmed <= 7.8e-6, below the smallest of the 25 boundary pmed
% (3.78e-5). (ii) E4 (results/e4/e4_tests.csv, test == critic): Spearman
% correlation of T between the two test seeds across the 30 universes of a generator, -0.03
% (MV bootstrap) to 0.30 (CCC-GARCH-t; 0.3046), as in tab:manual-e4-reliability. The architecture ablation (results/e5arch) is post hoc;
% the pre-specified verdicts are those of the MLP critic. The CNN was chosen after a pilot on four
% E2 designs at n = 500 (docs/DESIGN.md, change log 2026-11-16). Integration pass (2026-11-16):
% tab:e1-sim-se (0.480 vs 0.060), tab:e3-rv-baseline (0.69-0.94), tab:e4-utility-agg (1.038,
% 1.004); classifier checkpoints from app:impl-nets.
% =============================================================================================

The question in the title has a qualified answer. For the five univariate generators and the
S\&P~500 panel we studied, synthetic data from GBM, Merton and Heston calibrated on 2018--2022
are not good enough: their windows are detectably less like the 2023--2025 market than the
2018--2022 history they were fitted to ($\pmed<0.001$, 0.005 and 0.017 for the anchored test;
\cref{tab:e3-verdicts}). Each verdict concerns a model together with its calibration rule.
Fitted to a period that contains the 2020 crash, these generators reproduce its volatility level,
whereas the median realised volatility of the real evaluation windows is 0.65 times the \TRAIN{}
standard deviation, and the witness of the critic shows that this excess volatility is mainly
what the test detects (\cref{sec:emp-witness}). Heston's rejection is
borderline against the boundary replications of the realistic design (\cref{sec:emp-verdicts}),
and it concerns Heston with the coefficient of variation of the variance bounded by
$\sqrt{1.5}$, the bound that keeps our daily discretisation free of zero-variance days
(\cref{sec:data-generators}). \GARCHt{}, the best of the three likelihood models by every
likelihood criterion, is not detectably worse than history ($\pmed=0.432$), but the protocol
cannot call it \GOOD: under the design that keeps the real volatility path of the calibration
period, the rule detects a volatility 20\% too high only with probability 0.45 [0.36, 0.55]
(\cref{tab:e3-power}). The stationary bootstrap, which replays the past, receives no verdict:
with the split we ran, its windows duplicate windows the anchored critic was trained on, and the
split by source window that removed this leakage in simulation (\cref{tab:e2-relative-boundary})
was not applied to it. For portfolio construction the answer comes from the task, not from the
test, and here from one evaluation period in which correlations fell. The generators that gave
the best minimum-variance portfolios, CCC-GARCH-$t$ and MV Student-$t$, were those whose
correlations departed most from a past whose average correlation (0.419) was well above that of
the evaluation period (0.268; \cref{tab:e4-correlation}); under rising correlations this
association could reverse. The generator most often rejected, MV Gaussian, gave minimum-variance
portfolios with \TEST{} variance ratios close to one (1.038 replica-wise, 1.004 with pooled
moments; \cref{tab:e4-utility-agg}). Within a generator the critic statistic was too noisy across
synthetic seeds for its near-zero correlation with the loss (\cref{tab:e4-link}) to be
informative: across the 30 universes of a generator its Spearman correlation between the two test
seeds lies between $-0.03$ and 0.30 (\cref{tab:manual-e4-reliability}).

What the protocol guarantees is limited but explicit. With independent windows the permutation
p-value of the \WC{} test has exact level for every critic, and the permutation version of the
anchored relative test is exact at its boundary $\Law(G)=\Law(\Rpast)$ for every critic and every
law of the new data (\cref{prop:exact-level,prop:anchored}); the simulations confirm both
(\cref{sec:sim-size,sec:sim-relative}; the anchored test rejected at 0.065 and 0.055 at nominal
5\% at its boundary). On dependent panels the time-block split removes the common-date bias of
ticker splits (\cref{prop:ticker-bias}), whose naive version rejected a true null in 0.990 to
1.000 of the placebo replications, though not the bias that a persistent factor can carry into
adjacent blocks (\cref{cor:block-bias}). The Driscoll--Kraay p-value is asymptotically valid for a fixed critic
whose evaluation data are independent of its training data (\cref{prop:hac}), a premise that our
split without a gap does not ensure; in the placebo it rejected at 0.085 when the synthetic panel
shared the factor structure of the real one (\cref{tab:e1-size-sim}), but not when the synthetic
tickers were independent (\cref{sec:concl-limitations}). The verdicts use these Driscoll--Kraay
p-values, so a \BAD{} verdict is a rejection whose error rate at the relative boundary is
estimated, not assumed (0.08 [0.04, 0.15] in the realistic design); a \GOOD{} verdict would rest
on a power premise that cannot be verified (\cref{prop:verdict}). Within this protocol the critic
is one calibrated and interpretable tool, not the most powerful one. The pooled
Kolmogorov--Smirnov test wins against deviations of the marginal law, and kernel tests are at
least as powerful against changes of GARCH persistence at matched variance
(\cref{sec:sim-power}). The case for the critic is that it is one function: it can be evaluated
window by window, it detected dependence-only deviations that a marginal test cannot see
(\cref{fig:e2dep-power}), it ordered the three likelihood models as the predictive log-score
does (\cref{tab:e3-agreement}), an agreement that a random ordering reaches with probability
1/6, it extends to multivariate windows, and it supports a relative null that is exact at its
boundary.

\subsection{Limitations}\label{sec:concl-limitations}

\paragraph{Verdicts depend on the split.}
Under non-stationarity a verdict is a statement about one pair of periods. With the calibration
period ending in 2021 instead of 2022, every parametric generator is \INCONCL{}, and in the realistic design
the rule rejects the boundary generator in 0.60 [0.39, 0.79] of replications, as often as it
rejects the 20\% alternative (\cref{tab:manual-e5-split}). In the 15 replications in which both
are rejected, however, the alternative is far more significant: its $\pmed$ lies below every one
of the 25 boundary values, and the size-adjusted power of \cref{def:power-check} is still 0.60 at
a threshold where the boundary size is 0.04 (\cref{tab:manual-theory-calibration}). The failure
at that split is one of calibration rather than of discrimination, although the size-adjusted
power also stays below $\bmin$. The \BAD{}
verdicts of Merton and Heston also disappear under some one-at-a-time changes of the pipeline,
in 3 and 5 of the 13 changes (\cref{tab:manual-e5-sensitivity}). The decision rule
\eqref{eq:verdict} checks power but not boundary size, so it would have issued uninterpretable
\BAD{} verdicts at the alternative split had the anchored test rejected (\cref{sec:robust}). A
calibration of the generators on a recent part of the calibration period, which would separate
the model from the volatility level of its calibration window, was not run.

\paragraph{Driscoll--Kraay inference can be anti-conservative.}
The normal-reference Driscoll--Kraay p-value relies on few evaluation dates (about 15 per half in
E3). When the synthetic tickers are drawn independently, as for the univariate generators, the
learned tests over-reject in the placebo even with the time-block split (0.270 for the absolute
\WC{} test, 0.155 for the anchored test; \cref{tab:e1-size-sim}), and fixed-$b$ references do not
repair it. In the multivariate study, with 123 windows per sample and persistent volatility, the
\WC{} test rejected a true null at 0.27 for CCC-GARCH-$t$ and factor GARCH-$t$
(\cref{tab:e4-tests}); per ticker, at $L=5$, it rejected at 0.12 (\cref{tab:e3-perticker-null}).
We therefore size-adjusted or calibrated these p-values against simulated nulls, which is only as
good as the simulated null. The interleaved blocks without a gap leave dependence between
training and evaluation windows that the theory does not control (\cref{lem:decoupling}).
\Cref{cor:block-bias,rem:small-T} give a model-based and heuristic explanation of the oversize. A
critic trained on dates adjacent to the evaluation blocks carries the persistent realised factor
path into them, and with independent synthetic tickers the conditional bias of $T$ and its
Driscoll--Kraay standard error depend on the same real path. In the placebo the mean of $T$ is
nearly the same in the two cells, but with independent synthetic tickers $T$ and its standard
error are negatively correlated (Pearson $-0.64$, against 0.22 when the synthetic panel has a
factor of its own; \cref{tab:manual-e1-sim-stat}), so positive values of $T$ come with small standard errors and
$\zhac$ is shifted: with the \WC{} test, 0.480 of the replications with a standard error at or
below the median reject, against 0.060 of the others (\cref{tab:e1-sim-se}). An embargo between
training and evaluation blocks, as in $hv$-block and purged cross-validation
\citep{racine_2000_consistent,lopezdeprado_2018_advances}, a contiguous split or a critic trained
on an independent copy of the panel would test this explanation; none was run.

\paragraph{Power under realistic non-stationarity.}
In a stationary reference design the verdict rule detects a 20\% volatility error in every
replication; with the real volatility path of 2018--2022 it does so in 0.45 of them. The median
standard error of the anchored statistic at the boundary is 0.245 in that design against 0.032 in
the stationary one (\cref{tab:e3-power}), plausibly because volatility regimes common to all
tickers make the few evaluation dates the binding constraint (\cref{rem:small-T}). On this panel
and split no generator could have been \GOOD, which was known before any generator was tested in
the full run (\cref{sec:emp-power}). The power check is also one-sided in an economic sense.
Its deviation, a volatility too high, points away from the calmer evaluation period, whereas a generator error
that points towards the new regime makes the generator more like the new data than history is,
and the anchored null does not penalise it.

\paragraph{The critic is not a likelihood estimator.}
Learned scores rank windows but track the exact log-likelihood ratio poorly (Spearman
correlations from 0.024 to 0.556; \cref{tab:e2-llr}), WGAN-GP critics are poor estimators of
$\W$ \citep{mallasto2019well}, and on real data the critic behaved mostly as a realised-volatility
statistic. On the seed-0 windows of the absolute test, a fixed log realised-volatility statistic
reaches between 0.69 and 0.94 of the critic's $\zhac$ and ranks the three likelihood models in the
same order (\cref{tab:e3-rv-baseline}); learning added some precision but no new direction. We did
not run this natural non-learned baseline inside the anchored rule with five seeds and the designs
D1 and D2, so whether the verdicts need a learned critic is not shown.
Where a likelihood exists it is the sharper tool, and the critic statistic should be read as a
test statistic and a witness, not as a distance or a likelihood.

\paragraph{The gap to the oracle.}
The Neyman--Pearson oracle rejects in at least 0.98 of the replications at every independent-window
alternative, whereas every window test stays near its null rate against the smaller jump
deviations (\cref{tab:e2-power}). Part of the gap is structural: with a 50/50 split and equal
sample sizes, a test based on the held-out mean difference of a fixed critic has at most half the
oracle's local noncentrality (\cref{prop:local-power}). The rest is the cost of learning: against
a 5\% volatility increase with 2000 windows the \WC{} test rejects at 0.58, MMD at 0.93. The
\LRC{} and C2ST classifiers selected their first checkpoint, after 50 steps, in almost every
simulated fit (\cref{app:impl-nets}), so their power in \cref{sec:sim} is that of nearly
untrained classifiers; longer training without early stopping was compared only in a pilot. The
power of the critic also depends on its training. A learning rate of $10^{-4}$ raised its power against
a 10\% volatility increase from 0.40 to 0.78 (\cref{tab:e5-simpower}), and the verdicts reported
here are those of the pre-specified multilayer-perceptron critic. With a time-pooled
convolutional critic, examined after the real-data results were known, the five real-data
decisions are unchanged, and in the 20 replications of the realistic design that it shares with
the perceptron its rule power and boundary size equal those of the perceptron (0.45 and 0.20 in
that subset; \cref{tab:e5-arch-real}). In simulation with independent windows it keeps the level and raises
the power against every deviation we studied, for instance from 0.10 to 0.98 against
Student-$t_{10}$ tails and from 0.10 to 0.92 when the synthetic windows lack the jumps of the
real ones (\cref{tab:e5-arch}), which suggests that much of the power gap is architectural,
although the network was chosen after a pilot on some of these designs and its size under the
dependent protocol was not measured. We did not compare with power-optimised deep-kernel tests
\citep{liu2020learning}.

\paragraph{Scope.}
The panel consists of the S\&P~500 constituents at the end of the evaluation period. Both
periods are therefore selected on information from the end of \TEST{}, a look-ahead that goes
beyond survivorship bias, and part of the volatility gap between the generators and the real
\TEST{} returns may be a selection effect; a replication should draw the universe from the
constituents at the end of \TRAIN{} (\cref{sec:data-panel}). One market and one pair of periods
are studied. Univariate verdicts say nothing
about cross-sectional dependence (\cref{ex:corr-blind}). The Markowitz conclusions rest on five
multivariate generators, two test seeds, 60 placebo replications per generator and 30 universes
that share every evaluation date; the utility differences are point estimates without inference
over dates, the portfolios use only the unconditional moments of the generators, and several
planned replications were cut for compute (\cref{app:impl}). The
generators are classical parametric models, without asymmetric GARCH specifications, and
resampling schemes; the generative networks that motivate much of the literature were not evaluated.

\paragraph{Future work.}
Five extensions follow directly. First, the decision rule should also check calibration at the
boundary, for instance by turning \BAD{} into \INCONCL{} when the Clopper--Pearson interval of the
boundary size lies above $\alpha$, by reporting the size-adjusted power of \cref{def:power-check}
beside the nominal one, or by calibrating $\pmed$ against the boundary replications as a Monte
Carlo test (\cref{rem:calibrated-bad}), which is valid only insofar as the boundary design
reproduces the real configuration; a \GOOD{} verdict with a formal error guarantee could come from the
equivalence version of \cref{rem:equivalence}. Second, the dependent-data protocol needs a gap
between training and evaluation blocks, so that \cref{lem:decoupling} applies, together with
small-sample corrections calibrated on placebos that share the dependence of the generator under
test. Third, the rule needs more evaluation dates, from overlapping windows (a step shorter than
$L$, with Driscoll--Kraay inference) or a longer evaluation period, and verdicts should be
aggregated over several split points in a rolling-origin design, so that they describe a
generator rather than one pair of periods. Fourth, the real-data comparison should include a
fixed realised-volatility statistic in the same anchored rule with D1 and D2, generators
recalibrated on a recent window (for instance the last 615 \TRAIN{} days), the stationary
bootstrap, the historical-simulation benchmark of practitioners, tested with the split by source
window (which would also give a real-data placebo at the boundary), per-ticker nulls simulated
from each generator rather than from one pooled \GARCHt{} null, and a horizon-matched likelihood
benchmark such as the predictive log-score of 20-day returns. Fifth, power-optimised critics, architectures that pool over time and
pre-registered optimiser settings should be compared on the designs of \cref{sec:sim}, with the
size of each measured under the dependent protocol, and the protocol applied to neural generators
and to multivariate panels, where the use-based check of \cref{sec:markowitz} should accompany
every distributional verdict.
